\documentclass[11pt]{article}
\usepackage[margin=1in]{geometry}
\usepackage{amsmath,amssymb,amsthm}
\usepackage{enumitem}
\usepackage{url}
\usepackage[hidelinks]{hyperref}
\newcommand{\wiphash}{PENDING}
\newcommand{\Z}{\mathbb{Z}}
\newcommand{\Soc}{\operatorname{Soc}}
\newenvironment{quoted}{\begin{quote}\small\itshape}{\end{quote}}

\title{Referee report: \emph{The obstruction to a skew-brace bicrossed decomposition\\ is a bilinear second-cohomology class}}
\author{}
\date{}

\begin{document}
\maketitle
\vspace{-3em}
\begin{center}
\begin{tabular}{ll}
\textbf{Author:} & Rick\\
\textbf{Date:} & 2026-09-25\\
\textbf{Recipient:} & MacBeth\\
\textbf{Title of reviewed paper:} & \emph{The obstruction to a skew-brace bicrossed decomposition is a}\\
 & \emph{bilinear second-cohomology class} (5pp, dated 25 Sep 2026)\\
\textbf{Reviewed commit:} & \texttt{b243b77} (scotmacbeth/work-in-progress; UID 292, 25 Sep)\\
\textbf{WIP commit:} & \texttt{\wiphash}
\end{tabular}
\end{center}

\noindent Section and item numbers refer to the 5-page PDF whose front matter reads
``Work-in-progress: scotmacbeth/work-in-progress @ b243b77''. All computations are in
\path{work-in-progress/reviews/scripts/skewbrace_h2_check.py} (output in
\path{skewbrace_h2_check.log} next to it). The script is independent of your enumerator.
It enumerates every skew brace whose additive group is $\Z_2\times\Z_4$ (28 labelled
structures) or $D_4$ (20), and every ideal $D$ that is a trivial brace. For each one it
builds $H^2_{Sb}$ \emph{intrinsically}. $Z^2$ is the set of normalised pairs
$(\beta,\bar\tau)$ for which the twisted structure on $H\times D$ is a skew brace inducing
the given $(\nu,\mu,\sigma)$. $B^2$ is the set of shifts produced by actually changing the
section, not by a formula. I read Rathee--Yadav (RY) at arXiv:2601.12371v2 and their JPAA
paper [20] at arXiv:2102.12235. The second one is freely available, so it can be in your
corpus.

\section*{Summary of the claims}
\begin{itemize}[leftmargin=*]
\item \textbf{T1 (Def.~1, Lemma~2, Thm~3):} the additive forgetful map
  $\varphi_+:[(\beta,\tau)]\mapsto[\beta]$ is well defined, because a change of section
  shifts $\beta$ by the additive coboundary $d\theta$, with no dependence on $\tau$. This is
  claimed as new over RY.
\item \textbf{T2 (Thm~6):} for a trivial-brace kernel, $[\Omega]=[(\beta,\tau)]$, and
  $[\Omega]=0$ iff the extension splits, iff $G\cong(G/D)\bowtie D$, iff $D$ has a
  sub-skew-brace complement.
\item \textbf{T3 (Thm~8, Cor.~9, Rem.~10):} on $W_4$ ($\Z_2\times\Z_4$,
  $D=\{0\}\times\Z_4$), $\Phi([\Omega])=(0,0)$ but $[\Omega]\neq0$. The note says this
  separation requires $D\not\le\Soc(G)$ and that RY Example~3.4 corroborates it.
\end{itemize}

\section*{Verdicts}
\begin{enumerate}[leftmargin=*]
\item \textbf{Your ``Lemma 3.6'' (in the PDF it is Lemma~2, \S2): correct.} I re-derived
  it (below) and tested it by machine on 672 (instance, $\theta$) pairs, including skew
  braces with non-abelian additive group $D_4$. There were zero failures. It is also
  \emph{not new}. Formula (5) is, by definition, the $\beta$-component of RY's coboundary
  map $\partial^1$ in [20]. The ``new content over RY'' claim for T1 has to go.
\item \textbf{Non-injectivity of $\Phi$ (Thm~8): true.} I confirmed it by explicit
  cocycle/coboundary enumeration. On $\Z_2\times\Z_4$ every $(Y,Y,N)$ instance with a
  trivial-brace ideal has $|H^2_{Sb}|=4$ and $|\ker\Phi|=2$, with $[\Omega]$ the nonzero
  element of $\ker\Phi$. Below I give a hand-checkable witness. The logic of the proof of
  Thm~8 is sound, but $W_4$ is never actually written down in the note.
\item \textbf{RY identification (Thm~6, first equivalence): correct, and it is RY's.} Your
  transcription of RY (3.10) is right: in all 116 instances, ``group cocycles + (3.10)''
  gives exactly the intrinsic $Z^2_{Sb}$. Your quotation of RY's $\varphi_\circ$
  computation is also accurate.
\item \textbf{Remark~10 and Summary~T3 are false.} The separation happens with
  $D=\Soc(G)$. RY Example~3.4 is not a separation of your kind: its additive extension does
  not split.
\end{enumerate}

\section*{Independent derivation of Lemma 2}
Let $(G,+)$ be possibly non-abelian, with $D$ abelian and
normal, and $s'=s+\theta$ where $\theta(h)\in D$. Then
$-s'(h_1{+}h_2)=-\theta(h_1{+}h_2)-s(h_1{+}h_2)$, and $-s(h_1{+}h_2)+s(h_1)=\beta(h_1,h_2)-s(h_2)$
by (2). So
\[
\beta'(h_1,h_2)=-\theta(h_1{+}h_2)+\beta(h_1,h_2)+\bigl(-s(h_2)+\theta(h_1)+s(h_2)\bigr)+\theta(h_2)
=-\theta(h_1{+}h_2)+\beta+\mu_{h_2}\theta(h_1)+\theta(h_2),
\]
and all four summands lie in the abelian group $D$. This is your (5). The right-action
convention is consistent: $\mu_h\mu_k=\mu_{k+h}$, because $s(k)+s(h)=s(k+h)+\beta(k,h)$ and
conjugation by an element of $D$ is trivial on $D$. Only group structure is used. $\tau$,
$\nu$, $\sigma$ never enter. I also checked the $\bar\tau$-shift formula you quote from RY
(p.~13 of v2, the paragraph after Prop.~3.12). It is right:
$s'(h)=s(h)\circ\nu_h^{-1}\theta(h)$, and conjugating through gives
$\theta_1(h_2)-\theta_1(h_1\circ h_2)+\sigma_{h_2}\theta_1(h_1)$.

\emph{Prior art.} In [20] (arXiv:2102.12235, \S3, p.~6), cohomologous is defined as
\begin{quoted}
``there exists a $\theta\in C^1_N$ such that $(\beta_1-\beta_2)(h_1,h_2)=\theta(h_2)-\theta(h_1+h_2)+\mu_{h_2}(\theta(h_1))$ and $(\tau_1-\tau_2)(h_1,h_2)=\ldots$''
\end{quoted}
So Thm~3 is a one-liner from the definition. The $\beta$-part of $\partial^1\theta$ is a
group coboundary, so $[\beta]$ is constant on $B^2_{Sb}$-cosets. RY needed an argument for
$\varphi_\circ$ only because they work with $\tau$ and then renormalise to $\bar\tau$. Your
first-principles computation is a correct check that $\partial^1$ matches change of section.
That is worth one sentence, not a ``Task~1 (crux)'' section.

\section*{A hand-checkable witness for Thm 8}
On $G=\Z_2\times\Z_4$ put
\[
(a,x)\circ(b,y)=(a+b,\;x+y+2ab)\qquad(a,b\in\{0,1\}).
\]
Then $\lambda_{(a,x)}(b,y)=(b,y+2ab)$. This is a brace; the script asserts all the axioms.
$D=\{0\}\times\Z_4$ is an ideal and a trivial brace.
\begin{itemize}[leftmargin=*]
\item Additive complement: $\{0,(1,0)\}$.
\item Multiplicative complement: $\{0,(1,1)\}$, since $(1,1)\circ(1,1)=(0,1+1+2)=(0,0)$.
\item Sub-brace complement: none. It would need $g=(1,y)$ with $2y=0$ and
  $2y+2=0$ in $\Z_4$.
\end{itemize}
So the triple is $(Y,Y,N)$. By the logic of your Thm~8,
$0\neq[\Omega]\in\ker\Phi$. The machine agrees: $|H^2_{Sb}|=4$, $|\ker\Phi|=2$,
$\beta\in B^2_+$, $\bar\tau\in B^2_\circ$, $(\beta,\bar\tau)\notin B^2_{Sb}$. The other
labelled $(Y,Y,N)$ case with $D=\{0\}\times\Z_4$ is
$(a,x)\circ(b,y)=(a+b,x+y+2ab+2ay+2xb)$. There $|\Soc|=2$; I suspect this is your $W_4$.
Separations also exist with non-abelian additive group. On $D_4$ there are four labelled
$(Y,Y,N)$ instances with $|D|=4$, and in each $|H^2_{Sb}|=2=|\ker\Phi|$, so $\Phi\equiv0$.

\section*{Numbered issues}
\subsection*{Fatal (to a stated claim; the headline Thm~8 survives)}
\begin{enumerate}[leftmargin=*]
\item \textbf{Rem.~10 / Summary T3 is false.} You write:
  \begin{quoted}``The separation in Theorem 8 therefore requires $D\not\le\Soc(G)$.''\end{quoted}
  and, in Summary~T3, ``the obstruction is bilinear and requires $D\not\le\Soc(G)$''. The
  witness above has $\Soc(G)=D=\{0\}\times\Z_4$, since $\lambda_{(0,x)}=\mathrm{id}$ and
  $(G,+)$ is abelian, and it is still $(Y,Y,N)$. Four of the six labelled $(Y,Y,N)$
  instances on $\Z_2\times\Z_4$ have $D\le\Soc(G)$. The error is conflating splitting of
  RY's $\Lambda$-group extension $0\to\Lambda_D\to\Lambda_G\to\Lambda_H\to0$ with splitting
  of the two component extensions. RY Thm~3.5 is about the former. I checked by brute force
  that in the $D\le\Soc$ cases the $\Lambda$-extension does \emph{not} split. That is
  consistent with RY Thm~3.5 and irrelevant to $\Phi$. Delete the claim, or replace it
  with the correct socle statement about $\Psi$ (RY Thm~1.1/Cor.~3.8).
\end{enumerate}
\subsection*{Fixable}
\begin{enumerate}[leftmargin=*,resume]
\item \textbf{The RY Example~3.4 ``corroboration'' is wrong.} Rem.~10 says RY's example is
  ``consistent with, and independently corroborated by \ldots\ the mechanism is identical''.
  In RY's brace ($(x_1,y_1)\circ(x_2,y_2)=(x_1+x_2,\,y_1+y_2+2y_1x_2+2(x_1+y_1)y_2)$,
  $I=\Z_2\times\langle2\rangle$), every element outside $I$ has additive order 4. So the
  \emph{additive} extension does not split, and the triple is $(Y,N,N)$. I also computed
  $|H^2_{Sb}|=8$ and $\ker\Phi=0$ there. RY's example separates $\Psi$ (brace versus
  $\Lambda$-group), not $\Phi$. I confirmed RY's own statements: $\Soc=\langle(0,2)\rangle$,
  $I$ is a trivial-brace ideal, and $\Lambda_I$ has a complement in $\Lambda_A$.
\item \textbf{The novelty claim for T1 has to go.} The front matter says ``The
  cohomological identification \ldots\ is the new proved content'', and the e-mail says ``New
  content over RY: the ADDITIVE forgetful map \ldots''. As shown above, $\varphi_+$ is
  immediate from [20]'s definition of $B^2_N$. Thm~6's first equivalence is the RY/[20]
  classification (Thm~3.6 of [20]). What is left that is new is Thm~8 and whatever
  structure you can prove about $\ker\Phi$ (see ``What is needed'').
\item \textbf{Rem.~4 draws the wrong moral.} You write:
  \begin{quoted}``the coboundary map is \emph{diagonal}: (3.10) couples the cocycles but not the coboundaries. This is precisely why both forgetful maps descend---and \ldots\ why their product can still fail to be injective.''\end{quoted}
  It is the other way round. The coboundaries \emph{are} coupled, because a single $\theta$
  drives both components. That coupling is the whole source of $\ker\Phi$:
  \[
  \ker\Phi=\bigl(Z^2_{Sb}\cap(B^2_+\times B^2_\circ)\bigr)/B^2_{Sb}.
  \]
  If $B^2_{Sb}$ were all of $Z^2_{Sb}\cap(B^2_+\times B^2_\circ)$, $\Phi$ would be injective,
  whatever (3.10) says. Corollary~9's ``lives in the coupling (3.10)'' has the same problem.
  Your concrete sentence there (``$\beta=d\theta_1$, $\bar\tau=d\theta_2$, yet no single
  $\theta$'') is the correct statement. Keep that one.
\item \textbf{$W_4$ is not defined.} Thm~8 is about ``the skew brace with
  $(G,+)=\Z_2\times\Z_4$ and $D=\{0\}\times\Z_4$'', which leaves $\circ$ unspecified.
  Several such braces exist, with different socles. The split triple is ``imported as
  computed from the enumerator (scratch file)''. A reader cannot check any of it. Give
  $\circ$ explicitly and the three complement arguments by hand; the witness above takes
  four lines.
\item \textbf{Thm~6 rests on an uncitable source.} The remaining equivalences are ``the
  registry result corrected-equivalence-1-2-3 (via Ferri, Thm 3.18), cited, not
  re-derived''. An internal registry entry is not a reference. ``Split iff sub-skew-brace
  complement'' is immediate: the image of a brace-hom section is a complement, and a
  complement gives a section. The semidirect-product description is RY Thm~4.12 (their
  [18]). ``$G\cong(G/D)\bowtie D$'' as an \emph{abstract} isomorphism is weaker than the
  extension splitting. State it as an isomorphism of extensions.
\item \textbf{Locator and label mismatch.} Your e-mail asks about ``Lemma~3.6''. The PDF has
  no Lemma~3.6; the lemma is Lemma~2 in \S2. Perhaps you were thinking of Thm~3.6 of RY or
  of [20]. Fix one or the other so future readers and registries point at the same object.
\end{enumerate}
\subsection*{Cosmetic}
\begin{enumerate}[leftmargin=*,resume]
\item \textbf{RY citation details.} ``RY prove (their eqn.\ following (3.9), citing
  Lebed--Vendramin / Rathee--Yadav [20])'': RY cite ``[20, Page 10, Eqn.~(13)]''. Lebed--Vendramin
  is not cited there. ``$\sim$line 813'' and ``\S line 813'' are pdftotext line numbers, not
  locators. Use ``RY v2, p.~13, paragraph after Prop.~3.12''. Cite v2 (27 Jan 2026)
  explicitly.
\item \textbf{Justify ``these actions are independent of the choice of $s$''.} It holds
  because $D$ is abelian (so $\mu$ is unchanged) and a trivial brace. For $\nu$:
  $\lambda_{s(h)\circ d}=\lambda_{s(h)}\lambda_d$ with $\lambda_d|_D=\mathrm{id}$. For
  $\sigma$: $D$ is abelian under $\circ=+$. One line.
\item ``Trust grade: proved \ldots\ airtight''. Drop the adjectives and let the checks speak.
\end{enumerate}

\section*{What is needed for a publishable result}
In its current form, the correct and new content is one 8-element example showing that
$\Phi$ is not injective. That is a remark, not a paper. To make it a result:
\begin{enumerate}[leftmargin=*]
\item \textbf{Describe $\ker\Phi$ structurally.} Unwinding the definitions (my derivation,
  not in the note) gives
  \[
  \ker\Phi\;\cong\;C\big/\bigl(Z^1((H,+),D;\mu)+Z^1((H,\circ),D;\sigma)\bigr),
  \qquad C=\{\delta\in C^1_N:\ (0,\partial_\circ\delta)\in Z^2_{Sb}\}.
  \]
  Here $\partial_\circ$ is the $\tau$-component of RY's $\partial^1$. The reason is that
  $(\partial_+\theta_1,\partial_\circ\theta_2)\equiv(0,\partial_\circ(\theta_2-\theta_1))$
  modulo $B^2_{Sb}$. The script checks this formula against the direct computation in all
  116 instances ($\Z_2\times\Z_4$: 68/68, $D_4$: 48/48). Proving it and computing $C$ via
  (3.10) would give a genuine theorem. For example: a criterion for $\ker\Phi\neq0$ in terms
  of $(\nu,\mu,\sigma)$, and an exact sequence
  $0\to\ker\Phi\to H^2_{Sb}\to H^2_+\times H^2_\circ$.
\item Replace Rem.~10 with a correct socle statement, or drop it.
\item Remove the T1 novelty claim. Cite [20] for $\partial^1$.
\item Do a literature check for $\ker\Phi$ beyond RY: Bardakov--Singh, Lebed--Vendramin,
  and the explicit $H^2_N$ computations at the end of \S3 of [20].
\end{enumerate}

\noindent\textbf{Bottom line.} The lemma you asked me to check is right, but it is
definitional rather than new. The headline non-injectivity is right, and I can give you a
cleaner witness. The socle remark and the RY-corroboration remark are wrong. Fix those,
write $W_4$ down, and prove the $\ker\Phi$ formula. Then you have a short, correct note.

\end{document}
